\section{The City and the Republic}

\textbf{Georges Dumezil} famously showed that the structure of Indo-European societies was based on a trifunctional division: priests, warriors, and workers. Obviously, this corresponds perfectly with the caste structure that both \textbf{Rene Guenon} and \textbf{Julius Evola} regarded as one of the marks of a Traditional society. Conversely, one of the marks of the modern world is its departure from the trifunctional model with the result that hardly anyone knows what his ``natural'' place should be in the world. Hence, the following question must come to mind: How would a Traditional elite, should it arise again in the West, work to restore such a societal arrangement?

Clearly, the situation now is different from what it was at its origins. Trifunctionality developed organically in the original Indo-European groups. Men of different types and temperaments gravitated naturally to their proper positions. For example, a man became a priest because of his inner spiritual qualities and then arose to his place in the hierarchy. To be more explicit, this is the direct opposite of a ``social contract'', that is, a situation where individual men gathered together to debate and create the type of society that they voted and agreed on.

There would have been no need for such a discussion as the society would have preceded the individual. In The Ancient City, \textbf{Fustel de Coulanges} documents such societies in the Ancient Greek and Roman worlds. There was always a divine founder of the City, the lawgiver, who set up the constitution for the City. Piety was a primary virtue, since the City held together through time by the worship of the ancestral gods. To be impious was to be excluded from the City, which was tantamount to spiritual death.

In a subsequent stage of development, certain men, now known as ``philosophers'', began to reflect on the structure of the City. No longer content to blindly accept the status quo, they wondered about it: why is the City just? What makes a good priest, warrior, or artisan? What is the value of piety? And so on. What began in wonder resulted in complex discussions about the nature of reality, man, and society. This reached its high point in \textbf{Plato} and was codified by \textbf{Aristotle}. This was the complete philosophical reflection on the Greek City.

However, it also contained the seeds of destruction. The first question must be, ``Is the City just because it was constituted by the gods?'' This justice was simply assumed, but that is not acceptable to the philosopher. After many long dialogues, Plato's conclusion is that there is an idea of Justice that is independent of both man and gods. This is the basis for the charge of atheism against Socrates: if the ideas are primary, where is the need for the gods? What the people attribute to the gods, the philosopher attributes to the ideas. If the ideas are primary, then the gods are subject to the ideas and are therefore not free.

This puts the philosopher in an awkward position. He knows that piety is necessary for the continuance of the right order in the City. As a philosopher, he is rational and subjects himself to the just order. Unfortunately, the great majority of men are irrational and incapable of grasping the idea of justice. Moreover, they are weak and thus motivated by \emph{eros} and \emph{thymos} rather than \emph{nous}. For them, only piety towards the gods ensures the just order.

Unlike the numerous impious sophists of our day who preach and work toward revolution against the just order, Plato did not reject it. Instead, he created the just society in thought as the Republic. Obviously, it had to follow the trifunctional scheme as that represents the just order. The wise rule over the warriors and commoners. Obviously, the wise consist of those who know and are detached from material desires. The flaw in the Republic is that Plato does not know how to produce wise men apart from some strange and abhorrent practices.

In the next part, we will see how the Frankish bishops were not content to produce the just trifunctional society solely in thought, but made efforts to manifest it in the world of men and matter. Then we can wonder if a new elite would have any better success.


\flrightit{Posted on 2012-12-06 by Cologero}

\begin{center}* * *\end{center}

\begin{footnotesize}
\begin{sffamily}
\texttt{Dominion on 2012-12-06 at 00:32 said:}

In any Traditional order today, it's important to remember that while the natures may be the same (priest, warrior, worker), the world is not. The warrior, dealing with modern weaponry may find other ways to fight the active battle for Traditional principles. The worker, merchant and labourer, has access to technology and greater control over the means of production (which is certainly just, as we can find in the existence of the commons land in old Europe). What is necessary is that social institutions are such that each person can find his place according to his nature and true caste. The worker may prosper and develop while showing proper piety toward the higher principles, the warrior defends them, the Priest contemplates them and is part of the caste that creates the Intellect of the social order as a whole. The forms Tradition takes, as Evola says will be ``new and greater forms.''

\hfill

\texttt{Audrey on 2012-12-06 at 10:27 said:}

Quite interesting. Thanks again, Mr. Salvo. Just wondering how many would actually read it, and reflect on this.

\hfill

\texttt{Audrey on 2012-12-06 at 10:31 said:}

Fyi, just wondering if anyone has thoughts on the perversions of warrior caste, failure of the law, and technology making the worker lazy.

\hfill

\texttt{Andrew on 2012-12-06 at 16:39 said:}

I do. I think advancement of weapons technology is actually one of the main perverters of warrior culture. As Evola said, over time the warrior became a soldier, the soldier a serviceman, ``a sort of armed bureaucrat'' and now the serviceman is a geek behind a computer screen piloting drones. A real fight is face-to-face. The first perversion was projectile weapons.

The law has no more authority because nations, or rather States no longer recognize the legitimate spiritual authority real Law comes from.

And indeed, technology has made most persons into lazy cowards.

\hfill

\texttt{Mihai on 2012-12-07 at 03:51 said:}

To put it more briefly: any imbecile or couch potato can fire a gun and kill even the best of warriors from a safe (but within easy range) distance.

But it takes skill and discipline to wield a sword (like a Samurai warrior, for example) and a high spiritual constitution to use it as a support for the spiritual life.

So yes, firearms brought about ``egalitarism'' in war, where quantity replaced quality. It is no wonder that war in the modern times meant a global slaughterhouse, as the two world wars show.

\hfill

\texttt{Cologero on 2012-12-10 at 08:17 said:}

I believe, Audrey, that Plato addresses the warrior caste, or those with souls of brass ... aren't generals called ``the brass''? Their goal is power through the use of force and the acquisition of honors (e.g., praise, medals, citations, statues if they are lucky enough). That is their very nature, not necessarily a perversion. The normal and just order of things is for the warrior caste to be reined in and moderated by the intellectual/spiritual caste. The latter caste holds power through intellectual knowledge of the higher things and through persuasion, not brute force.

The perversion occurs then the reigns of power are held by the warrior caste --- called timocracy by Plato --- apart from the higher caste. This is an unjust order of things, hence a perversion. In Traditional terms, this is the ``degeneration of castes''.

\hfill

\texttt{Jason-Adam on 2012-12-10 at 19:09 said:}

Interesting how Plato considers Sparta a deviation from the perfect order whereas Evola, as most of the ``fascist'' right as opposed to the monarchist right, evokes Sparta as an ideal state... ... ... ... ..

Cologero, what exactly causes he degeneration of castes ? From my readings of Plato I cant quite see why the Republic beings to decline ? I see why once it begins to decline it cannot stop until it has become complete chaos but I do not know why a society begins to fall in the first place. The Evolian answer, based on Nietzsche, is that Europe fell because Christianity was somehow not right for it, but that explanation does not hold. So what causes the priestly caste to lose its hold over the masses ?

\hfill

\texttt{Cologero on 2012-12-10 at 21:30 said:}

Ironically, Nietzsche's and Evola's rejection of the spiritual basis of the last Traditional civilization of the West is itself a sign of the degeneration of castes.

Time is qualitative and unfolds in cycles. Why do you go to the seashore in Summer and ski in Winter? Because the conditions are right for it. Analogously, in different phases of the cosmic cycle, certain souls will be attracted to the material and psychic conditions of that phase. It will manifest through them. Others, however, a very few others who are awake and conscious, will also be attracted; they will be the harbingers of the end of one phase and the beginning of the next.

\hfill

\texttt{Jason-Adam on 2012-12-11 at 15:51 said:}

So do you mean that the cycle is fated from heaven and can not be stopped or begun by human choice ?

Also, do you see any significance in that the decline began in Europe and not somewhere else ? Did that mean European Tradition was somehow defective in that it could not build a stable society that lasted ?

\hfill

\texttt{Cologero on 2012-12-12 at 05:15 said:}

Guenon did not say the cycle is ``fated'', rather that it is the very nature of time. Can you conceive of space apart from the four directions? Of course not, you can just ``see'' it. So time has its cycle, although from a point we cannot (apart from mystical intuition) just ``see'' it as a whole. Since the owl of Minerva flies at dusk, at the end of a cycle we can understand it by looking back, but in a rational, not intuitive, way.

You are free to move around in space, at least within the limits of your powers. So your acts are free in time, not determined by the cycles. We have pointed that out in \textit{The Prophecies of Rene Guenon}\footnote{\url{http://www.gornahoor.net/?p=35}}. This does not mean a Tradition is intrinsically defective. Decline comes from forgetting Tradition, not from its alleged defects.


\hfill

\end{sffamily}
\end{footnotesize}

